\usepackage[margin=23mm,top=24mm,bottom=25mm,headheight=15pt]{geometry}
\usepackage[T1]{fontenc}
\usepackage{lmodern,microtype}
\usepackage{amsmath,amssymb,mathtools,bm,booktabs,array}
\usepackage{tikz,pgfplots}
\usetikzlibrary{arrows.meta,positioning,calc,decorations.pathmorphing}
\pgfplotsset{compat=1.18}
\usepackage[most]{tcolorbox}
\usepackage{fancyhdr,lastpage,hyperref,xcolor,enumitem}
\usepackage{needspace}
\definecolor{navy}{HTML}{202020}
\definecolor{teal}{HTML}{484848}
\definecolor{blue}{HTML}{383838}
\definecolor{orange}{HTML}{585858}
\definecolor{red}{HTML}{626262}
\definecolor{paper}{HTML}{F8F8F8}
\definecolor{lightteal}{HTML}{F8F8F8}
\definecolor{lightblue}{HTML}{F7F7F7}
\definecolor{lightorange}{HTML}{F8F8F8}
\definecolor{lightred}{HTML}{F7F7F7}
\hypersetup{colorlinks=true,linkcolor=blue,urlcolor=teal,
  pdftitle={The Nature of Noise: From Random Motion to Spectral Density, Thermal Equilibrium, and Circuit Analysis},
  pdfauthor={Mingyang Zhang},
  pdfsubject={Version 1.0. First public release: October 18, 2026.}}
\pagestyle{fancy}\fancyhf{}
\fancyhead[L]{\sffamily\small THE NATURE OF NOISE}
\fancyhead[R]{\sffamily\small ANALOG AND RF CIRCUITS}
\fancyfoot[L]{\footnotesize From probability to physical circuits}
\fancyfoot[R]{\footnotesize\thepage\,/\,\pageref{LastPage}}
\renewcommand{\headrulewidth}{0.4pt}
\setlength{\parindent}{0pt}\setlength{\parskip}{4pt}
\setlist[itemize]{leftmargin=2em,itemsep=2pt,topsep=2pt}
\setlist[enumerate]{leftmargin=2em,itemsep=4pt,topsep=3pt}
\newtcolorbox{concept}[1][]{enhanced,breakable,colback=lightblue,colframe=blue,
  boxrule=.65pt,arc=1.5mm,left=3mm,right=3mm,top=2mm,bottom=2mm,
  fonttitle=\bfseries\sffamily,title={Definition},#1}
\newtcolorbox{insight}[1][]{enhanced,breakable,colback=lightteal,colframe=teal,
  boxrule=.65pt,arc=1.5mm,left=3mm,right=3mm,top=2mm,bottom=2mm,
  fonttitle=\bfseries\sffamily,title={Physical interpretation},#1}
\newtcolorbox{example}[1][]{enhanced,breakable,colback=lightorange,colframe=orange,
  boxrule=.65pt,arc=1.5mm,left=3mm,right=3mm,top=2mm,bottom=2mm,
  fonttitle=\bfseries\sffamily,title={Worked example},#1}
\newtcolorbox{caution}[1][]{enhanced,breakable,colback=lightred,colframe=red,
  boxrule=.65pt,arc=1.5mm,left=3mm,right=3mm,top=2mm,bottom=2mm,
  fonttitle=\bfseries\sffamily,title={Convention and limitation},#1}
\newtcolorbox{derivation}[1][]{enhanced,breakable,colback=paper,colframe=navy!70,
  boxrule=.55pt,arc=1.5mm,left=3mm,right=3mm,top=2mm,bottom=2mm,
  fonttitle=\bfseries\sffamily,title={Derivation},#1}
\newcommand{\E}{\mathbb E}
\newcommand{\dd}{\mathrm d}
\newcommand{\Var}{\operatorname{Var}}
\newcommand{\Cov}{\operatorname{Cov}}
\newcommand{\rms}{\mathrm{rms}}
\newcommand{\enbw}{\mathrm{ENBW}}

% Keep related text on pages containing figures instead of creating sparse float pages.
\renewcommand{\topfraction}{0.9}
\renewcommand{\bottomfraction}{0.8}
\renewcommand{\textfraction}{0.08}
\renewcommand{\floatpagefraction}{0.8}
